\chapter{Addressing Complex Engineering Problems and Activities}

\section{Introduction}
Outcome-based engineering education assesses a final-year project against the
graduate-attribute framework of the International Engineering
Alliance~\cite{iea2021gapc}, under which the Washington Accord defines the
attributes of a Complex Engineering Problem (WP1 to WP7) and of Complex
Engineering Activities (EA1 to EA5). A problem qualifies as complex when it
demands in-depth knowledge (WP1) together with some or all of the remaining
attributes; an activity qualifies when its execution shows the corresponding
range of resources, interactions, and innovation. This chapter maps the
thesis onto those attributes, and each claim points at the section, table,
or figure of this report where the evidence sits.

\section{Complex Engineering Problems Associated with the Current Thesis}
\begin{itemize}
  \item \textbf{WP1: Depth of knowledge required.} Detecting sub-ten-pixel
        humans cannot be solved by applying a detector API with default
        settings. It required working knowledge of one-stage detection
        architectures at the layer level (Table~\ref{tab:layers}), of
        channel and spatial attention (Section~\ref{sec:cbam3}), of selective
        state-space models and their scan mechanics (Section~\ref{sec:mamba}), and of
        the loss geometry that makes few-pixel boxes unstable under strict
        IoU (Sections~\ref{sec:loss3} and~\ref{sec:metrics}), a fundamentals-first analysis rather
        than tool use.
  \item \textbf{WP2: Range of conflicting requirements.} The core tension is
        technical and non-technical at once: recall on the smallest targets
        against a parameter and latency budget for airborne hardware, and
        the humanitarian value of the system against its surveillance
        potential (Chapter~V). The resolution is measured, not asserted: the
        P2 scale buys very-tiny recall at $+0.5$~M parameters and $+1$~ms
        (Table~\ref{tab:main}), and the recall-weighted $F_2$ metric encodes
        the ethical asymmetry directly into the evaluation.
  \item \textbf{WP3: Depth of analysis, no obvious solution.} Which of
        attention, resolution, or sequence modelling helps tiny-target
        detection had no a-priori answer; plausible reasoning supported all
        three. The additive ablation (Section~\ref{sec:protocol}) is the analytical
        instrument built to answer it, and the outcome, one strong
        component, one cheap one, and one null, was not predictable from
        the literature.
  \item \textbf{WP4: Familiarity of issues.} The dominant issues are
        infrequently encountered in standard detection practice: targets
        smaller than one detection cell (Figure~\ref{fig:strideprob}), a
        dataset whose instances are $99.6\%$ below 32~px
        (Figure~\ref{fig:sizedist}), and a training collapse specific to
        the state-space blocks at epoch 28 (Figure~\ref{fig:training}).
  \item \textbf{WP5: Beyond standards and codes.} No engineering standard
        prescribes how to evaluate or deploy a tiny-human detector for
        search and rescue. The evaluation protocol, the frozen-split and
        provenance rules, and the operational threshold policy
        (Sections~\ref{sec:metrics} and~\ref{sec:protocol}) had to be designed for the problem.
  \item \textbf{WP6: Stakeholder scope.} The stakeholders of a rescue
        detector, operators who act on its boxes, agencies bound by aviation
        and data-protection rules, and the people being imaged, have needs
        that pull in different directions; Chapter~V works through those
        obligations and the design choices that answer them.
  \item \textbf{WP7: Interdependence.} The system is a chain of coupled
        sub-problems: dataset properties drive the architecture
        (Section~\ref{sec:p2novel}), the architecture drives memory behaviour and hence
        the training configuration (Section~\ref{sec:setup}), and the evaluation design
        determines whether any of it can be attributed. The single-variable
        ablation exists precisely to cut through that interdependence.
\end{itemize}

\section{Complex Engineering Activities Associated with the Current Thesis}
\begin{itemize}
  \item \textbf{EA1: Range of resources.} The work managed diverse
        resources: a 10{,}215-image benchmark with frozen splits, pretrained
        weights, a single 16~GB consumer GPU shared across 54 hours of
        training (Table~\ref{tab:carbon}), open-source frameworks, and a
        metric pipeline producing per-run manifests, CSV data, and plots so
        that every reported number is regenerable.
  \item \textbf{EA2: Level of interactions.} Theoretical design met hardware
        limits directly. The stride-4 scale quadruples head activation
        memory, which forced the physical batch to 8 with gradient
        accumulation holding the effective batch at 16 (Section~\ref{sec:setup}), and
        the automatic optimiser choice diverged on that same architecture,
        which forced the pinned AdamW configuration; resolving these
        interactions without breaking comparability was the day-to-day
        engineering of the thesis.
  \item \textbf{EA3: Innovation.} The work uses established principles in a
        way the source literature does not: attention added by substitution
        at negative parameter cost rather than as a bolt-on, state-space
        blocks placed surgically in the neck of an unchanged CNN detector so
        their effect is attributable (published Mamba detectors replace the
        backbone~\cite{wang2025mambayolo}), and an evaluation that elevates
        per-size recall and $F_2$ to first-class metrics because the task
        demands them.
  \item \textbf{EA4: Consequences for society and the environment.} The
        activity was run with its externalities measured: a carbon tracker
        and power sampler on every run (9.0~kg CO$_2$ total,
        Table~\ref{tab:carbon}), an ethical position on dual use and
        training-data dignity (Chapter~V), and a deliverable whose small
        size lowers the energy cost of every future deployment.
  \item \textbf{EA5: Familiarity, extension beyond prior experience.} The
        project extended well past coursework: reading and adapting current
        detection, attention, and state-space literature; diagnosing a
        silent model-rebuild fault that had invalidated an earlier line of
        results; and building failure-tolerant training infrastructure
        (smoke tests, resume, out-of-memory ladders) for a load-shedding
        environment, none of which is met in standard undergraduate
        laboratory work.
\end{itemize}

\section{Conclusion}
Measured against the Washington Accord attributes~\cite{iea2021gapc}, the
thesis presents a complex engineering problem, in-depth knowledge applied to
an unfamiliar, standard-less task with conflicting requirements and coupled
sub-problems, and its execution shows the corresponding complex activities,
from resource and interaction management to innovation and measured
consequences. The mapping is not decorative: each attribute above is
anchored to a numbered table, figure, or section of this report where the
claim can be checked.
